\section{Making Change}\label{greedy: sec: making change}\doHeader{Making Change}

Here's an example where, after the first greedy choice, the remaining problem
is just a smaller instance of the original problem.
So, the algorithm can simply recurse to solve that problem.
The example is simple, but nicely illustrates this particular schema for greedy algorithms.

\paragraph{Problem Definition.}
Making Change with U.S. coins.
\begin{mylist}
\item[\bf input:] an integer $N\ge 0$ (the desired total)
\item[\bf output:] A minimum-length multiset $\{c_1, c_2, \ldots, c_\ell\}$ such that each item $c_d$ is in $\{1,5,10,25\}$
  and the multiset sums to $N$, that is, $\sum_{i=1}^\ell c_i = N$.
\end{mylist}
(A \emph{multiset} is an unordered collection of items in which---in contrast to a set---any given item can occur more than once.)
Intuitively, the problem is to make change for a given value $N$, using as few (U.S.) coins as possible.
If item $c_i$ is, respectively, 1, 5, 10, or 25,
we think of it has being a penny, nickel, dime, or quarter.

\paragraph{Algorithm.}
There is a natural greedy algorithm---use the largest possible coin, then recurse:

\begin{myframed}
  \begin{steps}
  \item[{\textsf{rchangeUS}$(N)$:}]
    \step if $N = 0$: return $\emptyset$.  \comment{\ldots base case --- return the empty multiset}
    \step let $c' = \max\big\{c \in \{1, 5, 10, 25\} : c \le N\big\}$
    % [review 2026-09-27] fix: recursive call was named changeUS but the procedure is rchangeUS
    \step let $\{c_1, c_2, \ldots, c_{\ell'}\} = \textsf{rchangeUS}(N-c')$\comment{\ldots recursively make change for $N-c'$}
    \step return multiset $\{c', c_1, c_2, \ldots, c_{\ell'}\}$ \comment{\ldots return $\ell'+1$ coins of total value $N$}
  \end{steps}
\end{myframed}

Unwinding the tail recursion gives the following equivalent iterative implementation:
\begin{myframed}
  \begin{steps}
  \item[{\textsf{ichangeUS}$(N)$:}]
    \step while $N > 0$: 
    \step~~~let $c' = \max\big\{c \in \{1, 5, 10, 25\} : c \le N\big\}$
    \step~~~output $c'$
    \step~~~set $N = N-c'$
  \end{steps}
\end{myframed}

The run time is proportional to the number of coins output, which is $O(N)$.  

\paragraph{Correctness.}

To prove correctness, the main challenge is to show that the first greedy choice is correct---that is, there is always some correct solution that uses $c'$, the largest coin of value at most $N$.
Once we know that, we can commit to choosing coin $c'$ first,
and then (clearly?) the problem that remains is to make change optimally for the remaining total, $N-c'$.

We start with a useful observation:
\begin{observation}\label{greedy: obs: change}  For any $N$, every optimal way of making change has
  at most four pennies, at most one nickel,
  and at most two coins that are nickels or dimes.
\end{observation}
Each claim in the observation follows by an exchange argument.
(If there were more than four pennies,
the way of making change would not be optimal because five pennies could be exchanged for one nickel.
If there more than one nickel, two nickels could be exchanged for one dime.
And then, if there were more than two coins that are nickels or dimes,
there would have to be either three dimes (replaceable by a quarter and a nickel)
or two dimes and a nickel (replaceable by a quarter).)

Using the observation, we show that the first choice of the algorithm is correct:

\begin{lemma}\label{greedy: lemma: greedy change} 
  Fix any $N\ge 1$ and let $c' = \max\big\{c \in \{1, 5, 10, 25\} : c \le N\big\}$.
  Let $X^*$ be any optimal way of making change totaling $N$.
  Then $X^*$ contains a coin of value $c'$.
\end{lemma}
\begin{longFormProof}
  \step Consider any such $N$, $c'$, and $X^*$.
  \step\label{greedy: step: coin at most}
   All available values that are larger than $c'$ exceed $N$, so no coin of value greater than $c'$ is in $X^*$.
  \step To finish, we show that there is a coin of value at least $c'$ in $X^*$.
  \begin{block}
    \step \underline{Case 1.} Suppose $c'=1$. That is, $N\le 4$.
   \step $X^*$ has total value $N\ge 1$, so $X^*$ must have a coin of value at least 1.
  \end{block}
  \begin{block}
    \step \underline{Case 2.} Suppose $c'=5$.  That is, $5 \le N \le 9$.
   \step By Observation~\ref{greedy: obs: change},  the total value of the pennies in $X^*$ is at most 4.
    \step But $N\ge 5$, so there is at least one other kind of coin in $X^*$, necessarily of value at least 5.
  \end{block}
  \begin{block}
    \step \underline{Case 3.} Suppose $c'=10$.  That is, $10 \le N \le 24$.
    \step By Observation~\ref{greedy: obs: change}, the total value of the pennies and nickels in $X^*$ is at most $4+5=9$.
    \step But $N\ge 10$, so there is at least one other kind of coin in $X^*$, necessarily of value at least 10.
  \end{block}
  \begin{block}
    \step \underline{Case 4.} Otherwise $c'=25$.  That is, $N\ge 25$.
    \step By Observation~\ref{greedy: obs: change} the total value of the pennies, nickels, and dimes in $X^*$ is at most $4+20=24$.
    \step But $N\ge 25$, so there is at least one other kind of coin in $X^*$, necessarily of value 25.
  \end{block}
  \step By the case analysis, there must be a coin of value at least $c'$ in $X^*$.
  \step By this and Step~\ref{greedy: step: coin at most}, there must be a coin of value $c'$ in $X^*$.
\end{longFormProof}

By the lemma, we know that the first step made by the greedy algorithm is correct for any $N$.
That is, the value $c'$ that it chooses in $\{1, 5, 10, 25\}$ is guaranteed to be in any optimal solution with total $N$,
so it can commit to adding $c'$ to its solution.
Next we prove carefully that to make change of value $N$ optimally
it suffices to add $c'$ to any optimal way of making change of total value $N-c'$:

\begin{lemma}\label{greedy: lemma: recursive change}
  Fix any $N\ge 1$ and let $c' = \max\big\{c \in \{1, 5, 10, 25\} : c \le N\big\}$.
  Let $\{c_1, c_2, \ldots, c_{\ell'}\}$ be any optimal way to make change of value $N-c'$.
  Then $\{c', c_1, c_2, \ldots, c_{\ell'}\}$ is an optimal way to make change of value $N$.
\end{lemma}
\begin{longFormProof}
  \step Consider any such $N$, $c'$, and $c_1, c_2, \ldots, c_{\ell'}$.
  \step The $\ell'+1$ coins in $c', c_1, c_2, \ldots, c_{\ell'}$ do total $N$,
  because the coins in $c_1, c_2,\ldots, c_{\ell'}$ total $N-c'$.
  \step To finish, we show that to make change for $N$, at least $\ell'+1$ coins are required.
  \step Let $X^*=\{c^*_1, c^*_2, \ldots, c^*_{\ell}\}$ be an optimal way of making change for $N$.
  \step By Lemma~\ref{greedy: lemma: greedy change}, $c'$ is one of those coins.
  \step So $X^*$ contains $\ell-1$ coins that make change for $N-c'$.
  \step Since $\ell'$ is the minimum number of coins to make change for $N-c'$, it must be that $\ell-1 \ge \ell'$.
  \step Equivalently, $\ell \ge \ell'+1$---i.e., making change of total value $N$ requires at least $\ell'+1$ coins.
\end{longFormProof}

Finally, we show correctness.  
\begin{lemma}  
  For all $N\ge 0$, \textsf{rchangeUS}$(N)$ is correct---that is, it
  makes change totaling $N$ using the fewest possible coins.
\end{lemma}
\begin{longFormProof}
  \step The proof is by induction on $N$.
  \step[greedy: change base case]
  For base case $N=0$, \textsf{rchangeUS}(0) uses no coins, which is clearly the minimum possible.
  \begin{block}[greedy: change induction]
    \step Consider any $N\ge 1$, and assume that, for all $N'<N$, \textsf{rchangeUS}$(N')$ is correct.
    \step Consider the execution of \textsf{rchangeUS}$(N)$.  As in Lines 2 and 3 of the algorithm,
    \step define $c' = \max\{c\in\{1,5,10,25\} : c \le N\}$ and $\{c_1, c_2, \ldots, c_{\ell'}\}= \textsf{rchangeUS}(N-c')$.
    \step By the inductive assumption, $\{c_1, c_2, \ldots, c_{\ell'}\}$ is an optimal way to make change totaling $N-c'$.
    \step So, by Lemma~\ref{greedy: lemma: recursive change}, $\{c',c_1,\ldots, c_{\ell'}\}$ is an optimal way to make change totaling $N$.
    \step But $\{c',c_1,\ldots, c_{\ell'}\}$ is also the output of $\textsf{rchangeUS}(N)$.
    \step So $\textsf{rchangeUS}(N)$ outputs an optimal way to make change totaling $N$.
  \end{block}
  \step By Block~\ref{greedy: change induction}, for all $N\ge 1$,
  \textsf{rchangeUS}$(N)$ is correct if \textsf{rchangeUS}$(N')$ is correct for all $N'< N$.
  \step By Step~\ref{greedy: change base case}, \textsf{rchangeUS}(0) is correct.
  \step By the preceding two steps and induction, for all $N\ge 0$, \textsf{rchangeUS}$(N)$ is correct.
\end{longFormProof}

The theorem below summarizes what we've figured out.
\begin{theorem}
  There is an $O(N)$-time algorithm for \prob{Making Change with U.S. Coins}.
\end{theorem}

\subsection{Exercises}

\exercise \emph{(non-greedy denominations)}
The proof above relies heavily on the particular denominations $\{1,5,10,25\}$.
For each $d\in\{6,7,8,9,11\}$, do the following.
Consider using instead the set of denominations $\{1, 5, d\}$.
Figure out whether,
for all integers $N$, the greedy algorithm (with available denominations $\{1, 5, d\}$)
always makes change for $N$ using a minimum possible number of coins.
For each value of $d$, explain clearly why your answer is correct.

\paragraph{External exercise}
\begin{itemize}
\item  \JE Chapter 4, Exercise 14
\end{itemize}

%%% Local Variables:
%%% mode: latex
%%% TeX-master: "greedy_algorithms"
%%% End:
